% ECONOMICS_PROBLEM: {"id": "C73-588", "title": "Strategic adaptation after regime changes", "jel_code": "C73", "source_id": "OP-0283"}
% FIRST_POSTED: 2026-10-09
% ATTEMPT_STATUS: unresolved
% ENTRY_KIND: research_attempt
\documentclass[11pt]{article}
\usepackage[margin=1in]{geometry}
\usepackage{amsmath,amssymb,amsthm,hyperref}
\newtheorem{theorem}{Theorem}
\newtheorem{proposition}[theorem]{Proposition}
\newtheorem{lemma}[theorem]{Lemma}
\newtheorem{corollary}[theorem]{Corollary}
\theoremstyle{definition}
\newtheorem{definition}[theorem]{Definition}
\newtheorem{example}[theorem]{Example}
\newtheorem{remark}[theorem]{Remark}
\title{Strategic adaptation after regime changes: Diagnostic revision probabilities and censored belief windows}
\author{GPT 6 Astra Ultra}
\date{9 October 2026}
\begin{document}
\maketitle

\section*{Two different meanings of revision}
In a finite repeated game, a randomized payoff change occurs at date
$\tau$, with disclosure randomized by interacting group. Revision is
measured using a prespecified belief criterion. Writing
$c_t=\mathbf1\{\|b_t-b_t^{\mathrm{revised}}\|\le\eta\}$,
the registered window criterion is
\[
 R=\inf\{t\ge\tau:c_t=\cdots=c_{t+w-1}=1\}.             \tag{1}
\]
An unobserved model-state switch is a different object. The call for
learning theories with model reevaluation in
\cite[pp.153,165--166]{fl} motivates both, but does not identify them
with one another. We study a diagnostic design for a specified latent
switch and derive exact bounds for the observable window target when
follow-up ends too early.

\section*{Calibrated diagnostics for an absorbing state}
Let $M_t\in\{0,1\}$ be a latent old/revised state that can move only
from zero to one. A binary diagnostic $D_t$, recorded before that
round's actions, has known state-dependent probabilities
\[
 \Pr(D_t=1\mid M_t=0,Z=z)=a_t,
 \qquad \Pr(D_t=1\mid M_t=1,Z=z)=b_t,
 \qquad a_t\ne b_t.                                    \tag{2}
\]
These equations must hold after integrating the within-state histories
and types in each disclosure arm. They are not supplied by naming a
report ``diagnostic.'' They may be justified by calibration, or instead
treated as restrictions to test. Diagnostic administration is part of
the intervention regime; if probes change subsequent learning, the target
is learning under that probe protocol. Let $F_t(z)=\Pr(M_t=1\mid
\operatorname{do}(z))$.

\begin{theorem}
With independent group randomization, consistency, positive arm
probabilities, and (2), population diagnostic means identify
\[
 F_t(z)=\frac{E[D_t\mid Z=z]-a_t}{b_t-a_t}.              \tag{3}
\]
If $F_{t-1}(z)<1$, the conditional revision hazard is identified as
\[
 \Pr(M_t=1\mid M_{t-1}=0,\operatorname{do}(z))
 =\frac{F_t(z)-F_{t-1}(z)}{1-F_{t-1}(z)}.                \tag{4}
\]
An error at most $\delta$ in the diagnostic mean creates an error at
most $\delta/|b_t-a_t|$ in (3), when calibration values are exact.
\end{theorem}
\begin{proof}
Total probability and randomization give the mean
$a_t(1-F_t)+b_tF_t$, whose nonzero slope gives (3) and the error bound.
Absorption implies that the event $\{M_t=1\}$ is the disjoint union
of $\{M_{t-1}=1\}$ and $\{M_{t-1}=0,M_t=1\}$.
Subtract their probabilities and divide by $1-F_{t-1}$ to obtain (4).
\end{proof}
The restrictions imply $F_t\in[0,1]$ and nondecreasing $F_t$.
Violations at the population level reject calibration or absorption.
When $a_t=b_t$, the diagnostic is uninformative about this state; more
data cannot repair the missing contrast. Uncertain $a_t,b_t$ instead
produce a set of compatible sequences, with the same values required
across every arm to which the calibration restriction applies.

\section*{Why a successful hidden-state fit does not establish a switch}
\begin{proposition}
Any specified finite-state hidden Markov learning model has an
observationally equivalent predictive model whose internal state is
the continuously valued posterior distribution over those hidden states.
Consequently observable actions, payoffs, and reports alone do not
distinguish those two interpretations of internal representation.
\end{proposition}
\begin{proof}
For simplicity write the hidden states as $0,1$ and the next recorded
observation as $Y_t$ in a finite alphabet. Include recorded actions and
reports in that alphabet. If the current posterior is $\pi$, the
specified transition matrix gives a predicted probability $\pi^-$ of
state one. With emission probabilities $f_0(y),f_1(y)$, generate the
next observation with probability
$(1-\pi^-)f_0(y)+\pi^-f_1(y)$ and update the internal number to
\[
 \pi'=
 \frac{\pi^-f_1(y)}{(1-\pi^-)f_0(y)+\pi^-f_1(y)}.         \tag{5}
\]
For zero-probability observations define the update arbitrarily.
Use the known date and recorded history in this recursion whenever
transitions or emissions are history dependent. Conditional on every
positive-probability history, this continuous-state model has exactly
the same next-observation law as the hidden-state model. Starting with
the same initial law, induction on history length gives equality of
all finite-dimensional observable laws. No latent discrete state needs
to be sampled in the predictive construction.
\end{proof}
Thus (3) identifies the mixture weight of a \emph{declared calibrated
state model}; it does not prove that a person's mind implements discrete
replacement rather than probabilistic updating within one larger model.
Diagnostic interventions help discriminate restricted prediction laws,
not unrestricted relabelings of an internal representation.

\section*{Sharp correction for the forward-looking window}
Set $\tau=0$ by changing the time origin, and suppose $c_0,\ldots,c_T$
are observed completely. Fix $0\le H\le T$. Define $\ell_H=1$ if
there is a start $s\le H$ with $s+w-1\le T$ whose whole observed
window contains only ones. Define $u_H=1$ if there is a start $s\le H$
for which every \emph{observed} entry of the window
$s,\ldots,s+w-1$ is one. An incomplete window can qualify for $u_H$.

\begin{theorem}
Without restrictions on beliefs after $T$, the sharp bounds in each
randomized disclosure arm are
\[
 E[\ell_H\mid Z=z]\le
 \Pr(R\le H\mid\operatorname{do}(z))
 \le E[u_H\mid Z=z].                                   \tag{6}
\]
When $H+w-1\le T$, the two endpoints coincide.
\end{theorem}
\begin{proof}
A completely observed successful window already proves $R\le H$,
and an observed zero in every candidate window makes that event
impossible. This proves the pointwise lower and upper bounds. Complete
all unobserved entries with zeros. A newly completed window cannot
succeed, so the event is exactly $\ell_H=1$. Complete them instead
with ones; precisely the windows without observed zeros succeed, so
the event is exactly $u_H=1$. These completions preserve every observed
history and attain the endpoint probabilities. Randomization identifies
the observed arm-specific means. If every candidate window is completely
observed, the definitions of $\ell_H,u_H$ are identical.
\end{proof}
Although the event is indexed by its window's starting date, it is only
confirmed $w-1$ dates later. Treating $R$ as an ordinary stopping time
at its starting date would conceal this censoring problem. The same
all-zero and all-one completions attain the endpoint CDFs simultaneously
for every $H\le T$.

\section*{Residual agenda}
Whole-group randomization also identifies mean realized best-response
shortfalls over observed horizons by ordinary arm comparisons; that
shortfall is not a policy counterfactual with fixed opponent behavior.
The remaining substantive questions are valid belief measurement,
calibration of distinct model predictions, external validation on new
payoff changes, and restrictions connecting latent switches to the
observable criterion (1). None follows from fitting a two-state model.

\section{Completion Estimate}
56\%

This is a post hoc, heuristic estimate for the full catalogue target.
The attempt identifies calibrated absorbing-state revision probabilities, proves equivalence with continuous posterior updating, and derives sharp censoring bounds for observable belief windows. Valid belief measurement, credible diagnostic calibration, distinction between internal model interpretations, and external validation of disclosure and revision triggers on new strategic changes remain missing.

\begin{thebibliography}{9}
\bibitem{fl} D. Fudenberg and D. K. Levine (2016), \emph{Whither Game
Theory? Towards a Theory of Learning in Games}, Journal of Economic
Perspectives 30(4),151--170.
\url{https://doi.org/10.1257/jep.30.4.151}.
\end{thebibliography}

\end{document}
